\documentclass[aps,pra,reprint,floatfix,superscriptaddress,nofootinbib]{revtex4-2}
\usepackage{amsmath,amssymb,booktabs}
\usepackage[T1]{fontenc}
\newcommand{\Ftwo}{\mathbb F_2}
\newcommand{\rank}{\operatorname{rank}}
\newcommand{\Span}{\operatorname{span}}

\begin{document}
\title{A seven-$T$ fixed measurement accessing 220 four-qubit Pauli directions}
\author{Xiaojian Zeng}
\email{xiaojianzeng@foxmail.com}
\affiliation{PBC School of Finance, Tsinghua University, Beijing 100190, China}

\begin{abstract}
The non-Clifford gate cost of a single fixed quantum-state measurement is a finite-resource question distinct from the number of available outcomes. We give an explicit Clifford+$T$ circuit on four input qubits and seven pure stabilizer ancillas whose complete, 2048-outcome measurement has an exact operator-span rank of 220. The circuit uses seven $T$ or $T^\dagger$ gates. Its rank follows from an exact Pauli-branch enumeration: 260 nonzero contracted readout operators occupy 220 distinct four-qubit Pauli directions. This exceeds the rank 208 attained by a particular previously reported seven-$T$ construction at the same input size and gate count. The remaining 36 blind directions make our measurement informationally incomplete. The example gives a stronger attainable partial-information benchmark; it does not determine the maximum seven-$T$ rank or the minimum gate count for complete four-qubit tomography.
\end{abstract}
\maketitle

\section{Introduction}

A fixed generalized measurement can extract information from an unknown quantum state without changing the circuit between shots. Single-setting tomography has been demonstrated using enlarged local Hilbert spaces and symmetric informationally complete readout \cite{Stricker2022}. For qubit circuits restricted to stabilizer ancillas and Clifford+$T$ gates, the non-Clifford resource cost is a different question. Lo Monaco \emph{et al.} \cite{LoMonaco2025} give general bounds and constructions for the rank of the measurement frame. Their Theorem~7 supplies an \emph{attainable} rank of 208 for four input qubits and seven $T$ gates; it is not an upper bound on what another circuit might achieve.

Here we exhibit one physical circuit at that same input size and $T$ count with rank 220. The improvement is in the dimension of the exactly accessible operator space, not in finite-shot accuracy, noise tolerance, Clifford depth, or complete tomography. In particular, the unresolved question of whether any circuit with at most seven $T$ gates can reach the full four-qubit rank 256 remains untouched by this construction.

\section{Measurement model and explicit circuit}

Let $D$ consist of four arbitrary input qubits and $A$ of seven ancillas prepared in $|0^7\rangle$. All 11 qubits are measured at the end in the computational basis; there is no intermediate readout, postselection, or randomized circuit choice. For a unitary $U$ and $b\in\{0,1\}^{11}$, the effective data effect is
\begin{equation}
 E_b=(I_D\otimes\langle0^7|)U^\dagger|b\rangle\langle b|U(I_D\otimes|0^7\rangle).
 \label{eq:effect}
\end{equation}
Thus $E_b\succeq0$ and $\sum_bE_b=I_D$. On the 256-dimensional real space of Hermitian data operators, the measurement frame satisfies
\begin{equation}
\begin{aligned}
 F&=\sum_b|E_b\rangle\!\rangle\langle\!\langle E_b|,\\
 \rank F&=\dim_{\mathbb R}\Span\{E_b\}.
\end{aligned}
\label{eq:frame}
\end{equation}

Number the data wires $0,\ldots,3$ and ancilla wires $4,\ldots,10$. Binary masks $x,z\in\{0,\ldots,2047\}$ use bit $i$ for wire $i$, least significant bit first. Define the Hermitian Pauli
\begin{equation}
\begin{aligned}
 W(x,z)&=i^{\operatorname{wt}(x\mathbin{\&}z)}
          \prod_{i=0}^{10}X_i^{x_i}Z_i^{z_i},\\
 \omega((x,z),(x',z'))&=x\!\cdot\!z'+z\!\cdot\!x'\pmod 2 .
\end{aligned}
 \label{eq:pauli}
\end{equation}
The seven rotations $R_j=\exp[-is_j\pi W(x_j,z_j)/8]$ act in increasing $j$ order and have the exact masks in Table~\ref{tab:rotations}. Each is a Clifford conjugate of $T$ or $T^\dagger$, up to an irrelevant global phase; $T^\dagger=S^\dagger T$ costs one non-Clifford gate in this convention.

\begin{table*}[t]
\caption{Seven chronological Pauli rotations. The sign $s_j$ selects $T$ or $T^\dagger$ up to Clifford factors and global phase. Masks use the convention of Eq.~\eqref{eq:pauli}.}
\label{tab:rotations}
\begin{ruledtabular}
\begin{tabular}{c r r c@{\qquad}c r r c}
$j$&$x_j$&$z_j$&$s_j$&$j$&$x_j$&$z_j$&$s_j$\\
\hline
1&837&805&$+1$&5&636&1340&$-1$\\
2&1194&801&$+1$&6&1562&1695&$+1$\\
3&120&1509&$+1$&7&1029&124&$-1$\\
4&38&1453&$-1$&&&&
\end{tabular}
\end{ruledtabular}
\end{table*}

Following the rotations, a single stabilizer readout measures the maximal commuting group generated by the eleven Pauli labels in Table~\ref{tab:readout}. The labels are mutually commuting and binary independent. A Clifford $C$ maps the computational $Z_i$ generators to these labels, up to signs, so the actual terminal circuit is
\begin{equation}
 U=C^\dagger R_7R_6\cdots R_1.
 \label{eq:circuit}
\end{equation}
The accompanying witness specifies a chronological 132-operation $H/S/\mathrm{CX}$ realization of $C$ and therefore removes any synthesis ambiguity. An explicit Clifford+$T$ expansion has 333 gates, exactly seven of them $T$ or $T^\dagger$. The Clifford count is not claimed minimal.

\begin{table*}[t]
\caption{Binary $(x,z)$ labels of the eleven independent, commuting readout generators. Generator signs only relabel outcomes and do not affect the operator span.}
\label{tab:readout}
\begin{ruledtabular}
\begin{tabular}{c r r@{\qquad}c r r@{\qquad}c r r}
$i$&$x$&$z$&$i$&$x$&$z$&$i$&$x$&$z$\\
\hline
0&1968&1249&4&985&1777&8&1851&2019\\
1&1901&652&5&1880&532&9&346&1556\\
2&539&1310&6&1880&880&10&1330&641\\
3&24&292&7&1285&1213&&&
\end{tabular}
\end{ruledtabular}
\end{table*}

\section{Exact rank certificate}

The group generated by Table~\ref{tab:readout} contains $2^{11}$ Hermitian Paulis $Q$ with eigenvalue-sign conventions fixed by $C$. The Fourier transformation between the projectors of this complete commuting readout and its group elements is invertible. Consequently the span of the physical effects in Eq.~\eqref{eq:effect} has the same dimension as the span of the contracted Heisenberg operators
\begin{equation}
 A_Q=\langle0^7|R_1^\dagger\cdots R_7^\dagger Q R_7\cdots R_1|0^7\rangle.
 \label{eq:contract}
\end{equation}
For a Hermitian Pauli $P$ and Pauli $Q$, conjugation by a rotation in Eq.~\eqref{eq:circuit} leaves $Q$ unchanged when $P$ and $Q$ commute. When they anticommute, it gives
\begin{equation}
 e^{is\pi P/8}Qe^{-is\pi P/8}=\frac{Q+isPQ}{\sqrt 2}.
 \label{eq:branch}
\end{equation}
The two terms have distinct binary labels and nonzero coefficients. This elementary branch rule is enough to compute the rank exactly; no floating-point rank threshold is used.

Write $a_j$ for the seven-bit ancillary-$X$ mask of rotation $j$. The masks are
\begin{equation}
\begin{aligned}
 (a_1,\ldots,a_7)&=(52,74,7,2,39,97,64),\\
 \rank_{\Ftwo}[a_1\ \cdots\ a_7]&=7.
\end{aligned}
 \label{eq:ancillary-basis}
\end{equation}
For a given $Q$, the projector $\langle0^7|\cdot|0^7\rangle$ annihilates every Pauli branch with a nonzero ancillary-$X$ mask. A branch selecting the product term at rotation indices $S\subseteq\{1,\ldots,7\}$ has final ancillary-$X$ mask $x_A(Q)+\sum_{j\in S}a_j$. Equation~\eqref{eq:ancillary-basis} implies that \emph{at most one} branch subset can survive: its selection vector is the unique solution of $[a_1\ \cdots\ a_7]e=x_A(Q)$ over $\Ftwo$. If a requested product branch is unavailable because its current Pauli commutes with the corresponding axis, then $A_Q=0$. Otherwise $A_Q$ is a nonzero real multiple of a single four-qubit Pauli. Hence cancellation between surviving branches is impossible.

We enumerate all $2^{11}$ readout-group masks and follow each unique branch in the reverse rotation order $j=7,\ldots,1$. The Pauli update on a selected product branch is $(x,z)\leftarrow(x,z)+(x_j,z_j)$ over $\Ftwo$; it is allowed exactly when the current symplectic pairing with $(x_j,z_j)$ is one. Unselected branches leave the label unchanged and have a nonzero coefficient whether they commute or anticommute. This finite exact calculation produces the counts in Table~\ref{tab:counts}. The accompanying code and JSON record every resulting label, including all 36 absent data labels. Because distinct four-qubit Paulis are real-linearly independent, the 220 distinct nonzero labels prove the claimed rank.

\begin{table}[t]
\caption{Exact enumeration for the circuit of Tables~\ref{tab:rotations} and \ref{tab:readout}. The 260 surviving Fourier rows need not yield distinct data operators.}
\label{tab:counts}
\begin{ruledtabular}
\begin{tabular}{l r}
Readout-group elements and physical outcomes&2048\\
Vanishing contracted Fourier rows&1788\\
Nonzero contracted Fourier rows&260\\
Accessible data-Pauli labels (rank)&\textbf{220}\\
Absent data-Pauli labels&36\\
\end{tabular}
\end{ruledtabular}
\end{table}

An independent state-vector construction applies the synthesized gate sequence to all 16 embedded data basis states in the exact field $\mathbb F_{1009}$, where the required images of $i$ and $\sqrt2$ exist. It checks the isometry and $\sum_bE_b=I_{16}$, obtains effect-matrix rank 220, and gives a nonzero 220-by-220 minor (determinant 955 modulo 1009). This supplies an independent lower-rank check; the unique-branch argument supplies the characteristic-zero upper bound. The modular calculation alone would not certify the upper bound.

OpenAI Codex (model configuration \texttt{gpt-5.6-sol}, ultra reasoning setting) substantively assisted the circuit search, exact-certificate code generation, and the computational checks. Separate AI-assisted implementations reconstructed the primary two-$T$ source example, enumerated all 2048 contracted operators, and evolved the explicit state matrix over the finite field. Their inputs, algorithms, and outputs are supplied for human scrutiny; the independent implementations are not a substitute for human scientific review.

\section{Comparison and limits}

For the same $n=4$ and $t=7$, the constructive value of Theorem~7 in Ref.~\cite{LoMonaco2025} is
\begin{equation}
 B(4,7)=2^{-3}(3^3-1)(3^2-1)^2=208.
 \label{eq:prior}
\end{equation}
That theorem presents 208 as achievable, while leaving the maximum open. The present rank 220 raises the demonstrated attainable rank by 12 independent data-Pauli directions relative to that particular construction. Rank is invariant under unitary Clifford relabeling of the input and invertible relabeling of outcomes, so the 220-dimensional span cannot be equivalent in rank to its 208-dimensional span. This comparison makes no claim of global priority beyond the cited work or of optimality among seven-$T$ circuits.

The 36 missing directions imply that some distinct four-qubit density operators give the same ideal outcome distribution. The exact rank says nothing about the smallest nonzero frame eigenvalue, statistical variance, gate noise, finite-shot reconstruction, or hardware connectivity. The construction uses seven ancillas and a long Clifford network; neither is optimized. A rank-256 measurement might still exist with seven $T$ gates and a different circuit or ancilla count, and this example does not resolve that possibility.

\section*{Data and software availability}
The exact witness, the branch-enumeration program and result record, and the independent explicit-circuit replay and result record accompany this draft in its reproducibility source packet. The data are finite integers and exact-field outputs; no experimental data were used. A persistent public archive and its citation are pending an author decision before submission.

\section*{Acknowledgments}
OpenAI Codex (model configuration \texttt{gpt-5.6-sol}, ultra reasoning setting) also assisted source comparison and manuscript drafting. The substantive research use and computational checks are described in Sec.~III. These records are available for human author assessment; no claim of completed human or external peer review is made.

\bibliography{references}
\end{document}
